\begin{abstract}
Data availability is foundational to reproducible computational science, yet existing decentralized storage primitives force researchers to choose between expensive, synchronous on-chain persistence and unverifiable off-chain pinning. Ethereum EIP-4844 introduced transient data blobs as an economical consensus mechanism, but its design was targeted almost exclusively at Layer-2 transaction rollups. In this paper, we explore a non-rollup application model: using EIP-4844 blobs as a temporary, verifiable publication layer for structured scientific and machine-learning research artifacts. We design and implement an end-to-end reproducible pipeline featuring RFC~8785 canonical record serialization, modulus-safe field element packing (\texttt{FE31}), a three-tiered commitment architecture combining blob versioned hashes, whole-dataset content hashes, and record-level Merkle proofs, and a dual-horizon content-addressed archival bridge. We empirically evaluate our system across five real-world and synthetic dataset classes (OpenAlex scholarly metadata, NOAA climate observations, PyTorch Transformer training logs, synthetic scaling traces, and incompressible controls) up to 50~MiB. Our results demonstrate a $27.8\times$ theoretical uncompressed cost reduction and up to $43.8\times$ measured compressed cost reduction over standard EVM calldata, with encoding throughput exceeding 170~MiB/s on commodity workstations. However, an empirical multi-tier retrieval study on Ethereum Sepolia reveals that public consensus Beacon endpoints exhibit severe instability and rate-limiting for historical sidecars. We demonstrate that while Ethereum blobs provide an economical, consensus-committed publication window, durable long-term scientific preservation strictly requires a decoupled, content-addressed archival bridge. Our full artifact, contracts, and evaluation harnesses are open source and independently reproducible.
\end{abstract}
